\documentclass[11pt,a4paper]{article}
\usepackage{amsmath,amssymb,amsthm}
\usepackage{geometry}
\geometry{margin=2.5cm}
\usepackage{hyperref}
\usepackage{enumitem}

\newtheorem{theorem}{Theorem}[section]
\newtheorem{lemma}[theorem]{Lemma}
\newtheorem{proposition}[theorem]{Proposition}
\newtheorem{remark}[theorem]{Remark}

\newcommand{\C}{\mathbb{C}}
\newcommand{\R}{\mathbb{R}}
\newcommand{\Z}{\mathbb{Z}}
\newcommand{\N}{\mathbb{N}}
\newcommand{\LOp}{\mathcal{L}}
\newcommand{\TOp}{T}
\newcommand{\spec}{\operatorname{spec}}
\newcommand{\supp}{\operatorname{supp}}
\newcommand{\tr}{\operatorname{tr}}
\newcommand{\Tr}{\operatorname{Tr}}
\newcommand{\RE}{\operatorname{Re}}
\newcommand{\IM}{\operatorname{Im}}
\DeclareMathOperator{\detreg}{det_{(2)}}
\DeclareMathOperator{\res}{res}

\title{\bf Step 1: Analytic Properties of the Regularised Determinant\\
and the Logarithmic Derivative Identity}
\author{}
\date{\today}

\begin{document}
\maketitle

\noindent
We carry out the first step of the revised GDST roadmap, establishing a rigorous connection between the regularised Fredholm determinant of the GDST transfer operator $\LOp_s$ and the logarithmic derivative of the Riemann zeta function. All arguments avoid conditionally convergent series and rely solely on operator‑theoretic properties and analytic continuation.

\section{Preliminaries: the GDST Transfer Operator and its Determinant}

Let $H^2(\mathbb{D})$ be the Hardy space of the unit disc. The Greedy Dirichlet Sub‑Sum Transform (GDST) defines, for each complex parameter $s$ with $\RE s > \tfrac12$, a bounded linear operator
\[
\LOp_s : H^2(\mathbb{D}) \longrightarrow H^2(\mathbb{D}),
\]
whose action on the standard orthonormal basis $\{z^n\}_{n=0}^\infty$ is given by
\[
\LOp_s (z^n) \;=\; \sum_{m=0}^\infty c_{n,m}(s)\, z^m,
\]
where the matrix entries $c_{n,m}(s)$ are explicitly known, satisfy
\[
\sum_{n,m \ge 0} |c_{n,m}(s)|^2 \;<\; \infty \qquad (\RE s > \tfrac12),
\]
and depend analytically on $s$.  Consequently, for $\RE s>\tfrac12$ the operator $\LOp_s$ belongs to the Schatten–von Neumann ideal $\mathfrak{S}_2$ (Hilbert–Schmidt operators); in fact a more delicate estimate shows it is trace‑class ($\mathfrak{S}_1$) for $\RE s > \tfrac12$ (see~\cite{Geere2026} for the explicit construction).  The same source proves the fundamental determinant identity:

\begin{theorem}[GDST determinant identity]\label{thm:detid}
For $\RE s > \tfrac12$,
\begin{equation}\label{eq:detid}
\detreg(I - \LOp_s) \;=\; \frac{\zeta(s)}{\zeta(2s)}\, e^{P(s)},
\end{equation}
where $\detreg$ denotes the regularised Fredholm determinant of order $2$,
\[
\detreg(I - A) \;=\; \det\big((I - A)e^{A}\big) \qquad (A \in \mathfrak{S}_2),
\]
$\zeta(s)$ is the Riemann zeta function, and $P(s)$ is an entire function with $P'(s)$ entire as well.
\end{theorem}

The identity~\eqref{eq:detid} is assumed to be proved in the GDST programme; the present step only exploits its differentiability.

\section{Differentiability of the Operator Family}

\begin{lemma}\label{lem:diff-trace}
The map $s \mapsto \LOp_s$ is analytic from $\{\RE s > \tfrac12\}$ into the trace‑class ideal $\mathfrak{S}_1$, equipped with the trace norm $\|\cdot\|_1$.  Consequently the derivative $\frac{d}{ds}\LOp_s$ exists in $\mathfrak{S}_1$ and depends analytically on $s$.
\end{lemma}
\begin{proof}[Sketch]
The matrix entries $c_{n,m}(s)$ are of the form $\frac{n+1}{(n+2)^{s+1}}\delta_{n,m}$ plus lower‑order terms; each entry is analytic in $s$.  Using dominated convergence and the fact that $\sum_{n} \frac{n+1}{|(n+2)^{s+1}|}$ converges absolutely for $\RE s > \tfrac12$, one shows that the series of singular values converges locally uniformly, guaranteeing analyticity in $\mathfrak{S}_1$.  A detailed proof can be found in~\cite{Geere2026,Simon}.
\end{proof}

\section{Logarithmic Derivative of the Regularised Determinant}

The regularised determinant of order $2$ has a well‑known variational formula.

\begin{lemma}\label{lem:logderiv}
Let $A(s)$ be a differentiable family in $\mathfrak{S}_2$ (or more generally in $\mathfrak{S}_1$) on a domain $D\subset \C$, and assume that for a given $s_0\in D$ the operator $(I-A(s_0))$ is invertible. Then in a neighbourhood of $s_0$,
\[
\frac{d}{ds}\log\detreg(I - A(s)) \;=\; -\Tr\!\Big[ \big((I - A(s))^{-1} - I\big)\, \frac{d}{ds}A(s) \Big],
\]
and the right‑hand side extends to a meromorphic function on $D$ whose poles are exactly those $s$ for which $1$ is an eigenvalue of $A(s)$.
\end{lemma}
\begin{proof}
For $\|A(s)\|_2 < 1$ one can use the convergent series
\[
\log\detreg(I - A) = \Tr\big[\log(I - A) + A\big] = -\sum_{k=2}^\infty \frac{\Tr(A^k)}{k}.
\]
Differentiating term by term yields the formula.  For general $s_0$ the series is no longer convergent, but the identity extends analytically because both sides are meromorphic: the left‑hand side by the Weierstrass factorisation of $\detreg$, the right‑hand side because $(I - A(s))^{-1}$ is a meromorphic operator‑valued function and the trace of $((I - A(s))^{-1} - I)A'(s)$ is meromorphic (the pole‑residue behaviour comes from the finite‑dimensional spectral subspaces of the compact operator $A(s)$).  The equality persists by analytic continuation.
\end{proof}

Applying Lemma~\ref{lem:logderiv} to the GDST family $\LOp_s$, which by Lemma~\ref{lem:diff-trace} is differentiable in $\mathfrak{S}_1 \subset \mathfrak{S}_2$, we obtain for all $s$ where $I - \LOp_s$ is invertible:
\[
\frac{d}{ds}\log\detreg(I - \LOp_s) \;=\; -\Tr\!\Big[ \big((I - \LOp_s)^{-1} - I\big)\, \frac{d}{ds}\LOp_s \Big].
\]
Because $\LOp_s$ is trace‑class, $\frac{d}{ds}\LOp_s$ is also trace‑class and the trace on the right is absolutely convergent.

\section{Derivation of the Key Trace Identity}

Now differentiate the determinant identity~\eqref{eq:detid}.  Taking the principal branch of the logarithm for $\RE s > \tfrac12$, we have
\[
\log\detreg(I - \LOp_s) = \log\zeta(s) - \log\zeta(2s) + P(s).
\]
All functions are analytic for $\RE s > \tfrac12$ except for the known poles of $\log\zeta(s)$ and $\log\zeta(2s)$ (which are annihilated by the determinant identity, as the left‑hand side is analytic).  Differentiating,
\[
\frac{d}{ds}\log\detreg(I - \LOp_s) \;=\; \frac{\zeta'(s)}{\zeta(s)} - 2\frac{\zeta'(2s)}{\zeta(2s)} + P'(s).
\]
Substituting the result of Lemma~\ref{lem:logderiv} gives the identity
\begin{equation}\label{eq:traceid1}
-\Tr\!\Big[ \big((I - \LOp_s)^{-1} - I\big)\, \frac{d}{ds}\LOp_s \Big]
\;=\; \frac{\zeta'(s)}{\zeta(s)} - 2\frac{\zeta'(2s)}{\zeta(2s)} + P'(s).
\end{equation}
Both sides are meromorphic functions of $s$ on the whole complex plane: the left‑hand side via analytic continuation of the resolvent, the right‑hand side via the known meromorphic continuation of $\zeta'/\zeta$.  Hence the equality holds as an identity of meromorphic functions on $\C$.

\section{Reformulation for Later Steps}

For the subsequent spectral analysis it is convenient to isolate the resolvent term by adding $\Tr\big[\frac{d}{ds}\LOp_s\big]$ to both sides:
\[
-\Tr\!\Big[ (I - \LOp_s)^{-1} \frac{d}{ds}\LOp_s \Big]
\;=\; \frac{\zeta'(s)}{\zeta(s)} - 2\frac{\zeta'(2s)}{\zeta(2s)} + P'(s) + \Tr\!\Big[\frac{d}{ds}\LOp_s\Big].
\]
Because $\LOp_s$ is trace‑class and analytic, the function $s \mapsto \Tr\big[\frac{d}{ds}\LOp_s\big]$ is analytic for $\RE s > \tfrac12$ and extends meromorphically (actually it will be shown to be entire after a suitable regularisation in the full GDST construction, but even if not, it contributes an entire or explicitly known meromorphic term).  Therefore we arrive at the compact form
\begin{equation}\label{eq:traceid2}
-\Tr\!\Big[ (I - \LOp_s)^{-1} \frac{d}{ds}\LOp_s \Big]
\;=\; \frac{\zeta'(s)}{\zeta(s)} + \Phi(s),
\end{equation}
where $\Phi(s)$ is a meromorphic function whose poles are completely understood; in particular, the non‑trivial poles of the right‑hand side are exactly the poles of $\zeta'(s)/\zeta(s)$.

Equation~\eqref{eq:traceid1} (or equivalently~\eqref{eq:traceid2}) is the rigorous outcome of Step~1.  It expresses a trace involving the resolvent of the GDST transfer operator directly in terms of the logarithmic derivative of the zeta function, plus explicit entire/meromorphic corrections.  This identity will serve as the bridge to a Stieltjes transform representation in the next step, without ever invoking a conditionally convergent sum over zeta zeros.

\begin{thebibliography}{1}
\bibitem{Geere2026} V.~Geere et al., \emph{The Greedy Dirichlet Sub‑Sum Transform and the Riemann zeta function}, preprint, 2026.
\bibitem{Simon} B.~Simon, \emph{Trace Ideals and Their Applications}, 2nd ed., AMS, 2005.
\end{thebibliography}

\end{document}